% ---------------------------------------------------------------------------
% Anexo: Declaración de uso de Inteligencia Artificial Generativa (IAG)
% Transcripción a LaTeX del formulario oficial UC3M (TFG, curso 2026)
%
% Requiere en el preámbulo:
%   \usepackage{array}
%   \usepackage{tabularx}
%   \usepackage{amssymb}
%   \usepackage{enumitem}
%   \usepackage{hyperref}
% ---------------------------------------------------------------------------

% --- Macros auxiliares (colocar en el preámbulo, no aquí) ------------------
% Casilla vacía y casilla marcada
\newcommand{\casilla}{\makebox[1.4em][l]{$\square$}}
\newcommand{\casillam}{%
  \makebox[1.4em][l]{\rlap{$\square$}\hspace{0.17em}\raisebox{0.15ex}{\scriptsize$\checkmark$}}}
% Ancho útil de una celda que abarca las dos columnas de la tabla
\newcommand{\anchopleno}{\dimexpr\linewidth-2\tabcolsep-2\arrayrulewidth\relax}
% Estilo de las filas de enunciado
\newcolumntype{Y}{>{\raggedright\arraybackslash}X}


\begin{center}
  \textbf{\large DECLARACIÓN DE USO DE INTELIGENCIA ARTIFICIAL GENERATIVA (IAG)\\
  EN EL TRABAJO DE FIN DE GRADO (TFG)}
\end{center}

\bigskip

\noindent\textbf{He usado IAG en mi TFG}


\smallskip

\noindent\begin{tabularx}{\linewidth}{|Y|Y|}
\hline
\casillam\textbf{SÍ} & \casilla\textbf{NO} \\
\hline
\end{tabularx}


\section*{Parte 1: declaración sobre comportamiento legal, ético y responsable}

Ten presente que el uso de IAG conlleva unos riesgos y puede generar una serie de
consecuencias académicas graves: la evaluación de tu TFG puede verse comprometida si
el uso de la IAG comporta la utilización de datos de carácter confidencial, materiales
protegidos por derechos de autoría, o datos de carácter personal, y se hace sin cumplir
las condiciones exigidas en cada caso (autorización de los interesados, autorización de
los titulares, seguimiento de las instrucciones de la Universidad).

\medskip

\noindent\begin{tabularx}{\linewidth}{|Y|Y|}
\hline
\multicolumn{2}{|p{\anchopleno}|}{\textbf{Pregunta}} \\
\hline
\multicolumn{2}{|p{\anchopleno}|}{%
  \begin{minipage}[t]{\anchopleno}\vspace{2pt}
  \begin{enumerate}[label=\arabic*., leftmargin=1.6em, nosep, topsep=0pt]
    \item En mi interacción con herramientas de IAG he facilitado \textbf{datos de
      carácter confidencial} contando siempre con la debida autorización de los
      interesados. La confidencialidad abarca toda información que una persona u
      organización desea proteger por razones legales, comerciales, de privacidad o
      estratégicas (como patentes o secretos comerciales).
  \end{enumerate}
  \vspace{2pt}\end{minipage}} \\
\hline
\casilla SÍ, he usado estos datos con la autorización de los interesados
  & \casillam NO, no he usado datos de carácter confidencial \\
\hline
\multicolumn{2}{|p{\anchopleno}|}{%
  \begin{minipage}[t]{\anchopleno}\vspace{2pt}
  \begin{enumerate}[label=\arabic*., leftmargin=1.6em, nosep, topsep=0pt, start=2]
    \item En mi interacción con herramientas de IAG he facilitado \textbf{materiales
      protegidos por derechos de autoría} contando siempre con la autorización de los
      respectivos titulares.
  \end{enumerate}
  \vspace{2pt}\end{minipage}} \\
\hline
\casilla SÍ, he usado estos materiales con autorización de los titulares de derechos de
  autor; o bien sin ella porque se ajustan a una de las excepciones o límites que permite
  la ley:
  \begin{itemize}[leftmargin=1.2em, nosep, topsep=2pt]
    \item obra en dominio público
    \item obra licenciada (licencias Creative Commons)
    \item uso de fragmentos con fines de investigación (derecho de cita)
  \end{itemize}
  & \casillam NO, no he usado materiales protegidos por derechos de autoría \\
\hline
\multicolumn{2}{|p{\anchopleno}|}{%
  \begin{minipage}[t]{\anchopleno}\vspace{2pt}
  \begin{enumerate}[label=\arabic*., leftmargin=1.6em, nosep, topsep=0pt, start=3]
    \item En mi interacción con herramientas de IAG he facilitado \textbf{datos de
      carácter personal} con la debida autorización de los interesados.
  \end{enumerate}
  \vspace{2pt}\end{minipage}} \\
\hline
\casilla SÍ, he usado estos datos con autorización de los interesados y conforme a las
  instrucciones contenidas en la
  \href{https://docs.google.com/document/d/1YghVxFwo8a1VqdfnvJiZzlmAkfYepbDmD4g3hTOW9ro/edit?tab=t.0}{guía aprobada por la Universidad}
  & \casillam NO, no he usado datos de carácter personal \\
\hline
\multicolumn{2}{|p{\anchopleno}|}{%
  \begin{minipage}[t]{\anchopleno}\vspace{2pt}
  \begin{enumerate}[label=\arabic*., leftmargin=1.6em, nosep, topsep=0pt, start=4]
    \item Mi utilización de la herramienta de IAG ha \textbf{respetado sus términos de
      uso}, así como los principios éticos esenciales, no orientándola de manera
      maliciosa a obtener un resultado inapropiado para el trabajo presentado, es decir,
      que produzca una impresión o conocimiento contrario a la realidad de los resultados
      obtenidos, que suplante mi propio trabajo o que pueda resultar en un perjuicio para
      las personas.
  \end{enumerate}
  \vspace{2pt}\end{minipage}} \\
\hline
\casillam SÍ & \casilla NO \\
\hline
\end{tabularx}

\section*{Parte 2: declaración de uso técnico}

Utiliza el siguiente modelo de declaración tantas veces como sea necesario, a fin de
reflejar todos los tipos de iteración que has tenido con herramientas de IAG. Incluye un
ejemplo por cada tipo de uso realizado.

\medskip

\noindent\textbf{Declaro haber hecho uso del sistema de IAG} (Claude, de Anthropic, modelos de la
familia Claude, versiones disponibles entre 2025 y 2026) \textbf{para:}

\subsection*{Documentación y redacción}

\begin{itemize}[leftmargin=1.4em]
  \item Soporte a la reflexión en relación con el desarrollo del trabajo: proceso
    iterativo de análisis de alternativas y enfoques utilizando la IAG.

  Se ha empleado la herramienta para someter el documento a una revisión crítica sistemática frente a la matriz de evaluación, \textbf{empleando el prompt} 
  \enquote{actúa como catedrático y miembro de un tribunal evaluador exigente; audita el documento adjunto frente a la matriz de evaluación}, 
  \textbf{teniendo como interacción} un proceso iterativo de contraste en el que se aceptaron parte de las observaciones y se rebatieron y descartaron otras.

  \item Revisión o reescritura de párrafos redactados previamente.

  Se ha empleado la herramienta para la revisión de redacción de fragmentos ya escritos con el fin de encontrar inconsistencias en el texto redactado y posibles faltas ortográficas.

  \item Búsqueda de información o respuesta a preguntas concretas.
  
  Se ha empleado la herramienta para realizar consultas puntuales sobre convenciones de citación, formalización de contrastes estadísticos y errores de compilación en \LaTeX.
\end{itemize}

\subsection*{Desarrollar contenido específico}

Se ha hecho uso de IAG como herramienta de soporte para el desarrollo del contenido
específico del TFG, incluyendo:

\begin{itemize}[leftmargin=1.4em]
  \item Asistencia en el desarrollo de líneas de código (programación).
  
  Depuración de errores en el código \LaTeX y revisión de fragmentos de código Python del análisis de datos.

  \item Tratamiento de datos: recogida, análisis, cruce de datos.

  Verificación de cálculos y contraste de decisiones metodológicas sobre el ajuste de distribuciones.
\end{itemize}



\section*{Parte 3: reflexión sobre utilidad}

% Valoración personal (formato libre) sobre fortalezas y debilidades identificadas en el
% uso de herramientas de IAG durante el desarrollo del trabajo: utilidad en el proceso de
% aprendizaje, en el desarrollo y en la extracción de conclusiones.

El uso de estas herramientas ha resultado especialmente útil en la fase de revisión crítica del documento,
al permitir un contraste sistemático frente a los criterios de evaluación que difícilmente se obtiene mediante
relectura propia. Su utilidad ha sido, en cambio, limitada en las decisiones metodológicas de fondo, ya que varias de las
observaciones recibidas partían de premisas incorrectas sobre el trabajo y era necesario rebatirlas y descartarlas, lo que
exigió mantener en todo momento el criterio propio como instancia última de decisión. La conclusión que se extrae es que
estas herramientas resultan valiosas como interlocutor crítico y como apoyo en tareas instrumentales, pero no sustituyen el
juicio técnico del autor sobre el dominio del problema.